% Standalone problem document, generated from the adjacent README.md.
% Compile with XeLaTeX. Mathematical content is maintained in Markdown.
\documentclass[11pt,a4paper]{article}
\usepackage[fontsize=10.5pt]{scrextend}
\usepackage[margin=22mm,headheight=15pt,headsep=9mm,footskip=11mm]{geometry}
\usepackage{fontspec}
\setmainfont{texgyrepagella}[Extension=.otf,UprightFont=*-regular,BoldFont=*-bold,ItalicFont=*-italic,BoldItalicFont=*-bolditalic]
\setsansfont{texgyreheros}[Extension=.otf,UprightFont=*-regular,BoldFont=*-bold,ItalicFont=*-italic,BoldItalicFont=*-bolditalic]
\setmonofont{texgyrecursor}[Extension=.otf,UprightFont=*-regular,BoldFont=*-bold,ItalicFont=*-italic,BoldItalicFont=*-bolditalic,Scale=MatchLowercase]
\usepackage{amsmath,amssymb}
\usepackage{unicode-math}
\setmathfont{texgyrepagella-math.otf}
\usepackage{xcolor,microtype,enumitem,needspace,fancyhdr}
\usepackage{bookmark}
\definecolor{ink}{HTML}{17344B}
\definecolor{accent}{HTML}{197D80}
\definecolor{muted}{HTML}{596773}
\hypersetup{colorlinks=true,linkcolor=accent,urlcolor=accent,pdftitle={MI-03 - Sharp
additive triangle constant for an odd number of
contractions},pdfauthor={Open Problems in Numerical Linear Algebra}}
\urlstyle{same}
\setlength{\parindent}{0pt}
\setlength{\parskip}{5pt plus 1pt minus 1pt}
\linespread{1.04}
\setlength{\emergencystretch}{3em}
\widowpenalty=10000
\clubpenalty=10000
\displaywidowpenalty=10000
\allowdisplaybreaks[1]
\setlist{nosep,leftmargin=1.5em}
\providecommand{\tightlist}{\setlength{\itemsep}{0pt}\setlength{\parskip}{0pt}}
\providecommand{\pandocbounded}[1]{#1}
\pagestyle{fancy}
\fancyhf{}
\fancyhead[L]{\sffamily\footnotesize\color{muted}OPEN PROBLEMS IN NUMERICAL LINEAR ALGEBRA}
\fancyhead[R]{\sffamily\bfseries\footnotesize\color{accent}MI-03}
\fancyfoot[L]{\parbox[b]{0.9\textwidth}{\sffamily\fontsize{8}{10}\selectfont\color{muted}Editorial ratings; a bounded search cannot certify that a problem remains unsolved.\\Literature check: 2026-09-14}}
\fancyfoot[R]{\sffamily\footnotesize\color{muted}\thepage}
\renewcommand{\headrulewidth}{0.3pt}
\renewcommand{\headrule}{\hbox to\headwidth{\color{accent}\leaders\hrule height \headrulewidth\hfill}}
\makeatletter
\renewcommand{\section}{\Needspace{5\baselineskip}\@startsection{section}{1}{0pt}{12pt plus 2pt minus 2pt}{4pt}{\sffamily\bfseries\large\color{ink}}}
\renewcommand{\subsection}{\Needspace{4\baselineskip}\@startsection{subsection}{2}{0pt}{9pt plus 1pt minus 1pt}{3pt}{\sffamily\bfseries\normalsize\color{ink}}}
\makeatother
\setcounter{secnumdepth}{0}
\begin{document}
{\sffamily\small\color{muted}Matrix inequalities and norms\par}
\vspace{3pt}
{\raggedright\hyphenpenalty=10000\exhyphenpenalty=10000\sffamily\bfseries\fontsize{21}{24}\selectfont\color{ink}Sharp
additive triangle constant for an odd number of contractions\par}
\vspace{7pt}
{\sffamily\small\textbf{Difficulty:} challenging\quad\textbf{Importance:} interesting
to specialist\par}
{\sffamily\small\bfseries\color{accent}Status: Lean verified\par}
\vspace{4pt}
{\color{accent}\hrule height 0.8pt}
\vspace{6pt}
\textbf{Rating rationale:} Sharpness for every odd summand count needs a
new extremal construction or obstruction; the payoff is concentrated in
operator triangle inequalities.

\subsection{Resolution --- 2026-09-11}\label{resolution-2026-09-11}

\textbf{Affirmative result by Matthew J. Colbrook}, Department of
Applied Mathematics and Theoretical Physics, University of Cambridge.
\textbf{Independent proof review: PASS.}

The sharp additive contraction constant is \(c_k=k/4\) for every
\(k\ge2\), including every odd summand count. Exact dimension-two
extremizers match the universal upper bound.

The exact target is resolved. The original statement and source evidence
are retained below; its former difficulty rating is historical.

\textbf{Primary manuscript:}
\href{https://github.com/ajt60gaibb/OpenProblemsInNLA/blob/main/matrix-inequalities-and-norms/MI-03/solution.pdf}{complete
proof PDF},
\href{https://github.com/ajt60gaibb/OpenProblemsInNLA/blob/main/matrix-inequalities-and-norms/MI-03/solution.tex}{standalone
TeX}, Theorem 1.1 and its proof;
\href{https://github.com/ajt60gaibb/OpenProblemsInNLA/blob/main/matrix-inequalities-and-norms/MI-03/solution.md}{authorship
and scope}. The
\href{https://github.com/ajt60gaibb/OpenProblemsInNLA/blob/main/references/colbrook-matrix-2026-09-11/verification/reviews/MI-03-review.md}{independent
review} checks the full original argument and records its hash. The
draft was AI-assisted; this is independent agent verification, not
external human peer review or formal certification.
\href{https://github.com/ajt60gaibb/OpenProblemsInNLA/blob/main/references/colbrook-matrix-2026-09-11/README.md}{Submission
record}.

\subsection{Lean proof and verification evidence ---
2026-09-14}\label{lean-proof-and-verification-evidence-2026-09-14}

\textbf{The complete original target is Lean verified.} The complete
original odd-summand conjecture is proved and, more strongly, the sharp
additive contraction constant is k/4 for every k≥2. The formalization
uses the actual complex Euclidean operator norm, positive square-root
matrix modulus and positive-semidefinite order, with genuine infima of
nonempty sets bounded below and exact dimension-two extremizers.

\href{https://github.com/sidneyholden1/OpenProblemsInNLA/tree/4602650e944c7221952554f495ff76c39c8b0708/matrix-inequalities-and-norms/MI-03/lean}{Immutable
formal proof}. The checked exports are \textit{NLA.\allowbreak{}MI03.\allowbreak{}upper\_\allowbreak{}bound},
\textit{NLA.\allowbreak{}MI03.\allowbreak{}sharpness}, \textit{NLA.\allowbreak{}MI03.\allowbreak{}sharp\_\allowbreak{}constant},
\textit{NLA.\allowbreak{}MI03.\allowbreak{}odd\_\allowbreak{}sharp\_\allowbreak{}constant}. Lean 4.33.1 uses pinned Mathlib
\textit{0df444a3} and LeanCert \textit{621a43d7};
\href{https://github.com/sidneyholden1/OpenProblemsInNLA/tree/codex/lean-mi03/matrix-inequalities-and-norms/MI-03/lean/README.md}{complete
pins and reproduction instructions} are retained. All four exported
transitive axiom closures contain only \textit{propext},
\textit{Classical.choice} and \textit{Quot.sound}.

\textbf{Formalization:} Sidney Holden, Center for Computational Biology,
Flatiron Institute, Simons Foundation, with OpenAI Codex assistance.
Matthew J. Colbrook retains mathematical authorship. Two independent
statement reviews preceded implementation, and two independent final
proof reviews passed. Kernel-only LeanCert and
\href{https://github.com/sidneyholden1/OpenProblemsInNLA/actions/runs/34862171380}{fresh
isolated Linux Comparator and kernel verification} passed, including
rejection controls.
\href{https://github.com/sidneyholden1/OpenProblemsInNLA/tree/codex/lean-mi03/matrix-inequalities-and-norms/MI-03/lean/verification/linux-2026-09-14/OPERATIONAL-REVIEW.md}{Original
artifact, exact checked inputs and independent operational audit} are
retained. This is AI-assisted formal verification, not external human
peer review or source-author endorsement.

\newpage

\subsection{Problem statement}\label{problem-statement}

For each integer \(k\ge2\), let \(c_k\) be the infimum of all \(c\ge0\)
such that

\[\left|\sum_{j=1}^k A_j\right|\preceq cI_n+\sum_{j=1}^k|A_j|\]

for every \(n\ge1\) and all \(A_1,\ldots,A_k\in\mathbb C^{n\times n}\)
satisfying \(\|A_j\|_2\le1\). Here \(|A|=(A^*A)^{1/2}\), \(\|\cdot\|_2\)
denotes the operator norm, and \(X\preceq Y\) means \(Y-X\) is positive
semidefinite.

Is \(c_k=k/4\) for every odd integer \(k\ge3\)?

\subsection{Why it matters}\label{why-it-matters}

The scalar absolute-value triangle inequality fails in matrix order. The
optimal additive correction quantifies that failure for uniformly
bounded summands, in a form usable in positive block-matrix estimates.

\subsection{References}\label{references}

\begin{enumerate}
\def\labelenumi{\arabic{enumi}.}
\tightlist
\item
  J.-C. Bourin and E.-Y. Lee, \emph{Diagonal and off-diagonal blocks of
  positive definite partitioned matrices}, arXiv:2307.02034v3 (15
  December 2023), Corollary 4.4 and the conjecture immediately following
  Remark 4.5. \href{https://arxiv.org/html/2307.02034}{Primary text}.
\item
  E.-Y. Lee, \emph{How to compare the absolute values of operator sums
  and the sums of absolute values?}, Operators and Matrices 6(3) (2012),
  613--619; §2 supplies related triangle inequalities.
  \href{https://files.ele-math.com/articles/oam-06-42.pdf}{Primary
  paper}, \href{https://doi.org/10.7153/oam-06-42}{DOI}.
\end{enumerate}

\subsection{Status check --- 2026-09-10}\label{status-check-2026-09-10}

Corollary 4.4 establishes \(c_k\le k/4\) for every \(k\) and sharpness
for even \(k\); the source separately conjectures sharpness for all odd
\(k>1\). Its latest version remains v3. Searches included
\textit{Bourin Lee contractions odd k constant k/4},
\textit{three contractions 3/4 sharp conjecture}, and
\textit{Bourin contractions sharp 2025 2026}. The 2026
symmetric-modulus papers concern different inequalities. No resolution
of this additive odd-summand problem was located.

\textbf{Audit update (2026-09-10):} Rechecked the conjecture after
Remark 4.5 in Bourin--Lee and searched for odd-summand sharpness
results. The proved even-summand statement is outside the target, which
already restricts to odd \(k\). This is a bounded literature check, not
a proof that no solution exists.
\end{document}
